\chapter{概览与物理背景}\label{chap:1}

\SourceParagraph{ch01-intro-001}{11}
本书介绍理论物理与数学交界处的一组主题，它们统称为“镜像对称”。镜像对称的概念于20世纪80年代末从超弦紧化的研究中逐渐形成，并在 Candelas、de la Ossa、Green 和 Parkes 于1991年发表的工作\BookCite{ref085}中首次得到精确表述。他们根据坚实的物理论证提出了一个猜想：五次 Calabi--Yau 流形上给定次数的有理曲线的数目，可以用另一个“镜像”Calabi--Yau 流形上的全纯三形式的周期来表示。沿着这一方向的发展还包括 Batyrev 对环面簇中超曲面给出的普遍镜像对称构造，以及 Givental 和 Lian、Liu、Yau 对 Candelas 等人所预测的瞬子数的证明。用物理语言来说，这些发展都涉及以 Calabi--Yau 空间为靶空间的 $N=2$ $\sigma$ 模型的 A 型与 B 型拓扑扭曲之间的关系，并且都与这些空间中的闭弦理论有关。本卷的前作\BookCite{ref242}以及文献\BookCite{ref101,ref458}对这些主题作了深入介绍。

\SourceParagraph{ch01-intro-002}{11}
20世纪90年代中期，在开弦理论的物理启发下，出现了两项更大胆的发展：Kontsevich 于1994年提出的同调镜像对称\BookCite{ref309}，以及 Strominger、Yau 和 Zaslow 于1996年提出的几何图景\BookCite{ref433}。这些思想使镜像对称超越了代数几何中的计数问题和物理中的二维 $\sigma$ 模型这两个较为专门的领域，形成了对这两个学科都有更广泛意义的图景。这两项发展及其所激发的研究，正是本书的主题。

\SourceParagraph{ch01-intro-003}{11}
我们首先回顾弦理论的一些一般物理背景，并对弦紧化、Dirichlet 膜、T 对偶以及其他物理概念作直观说明；这些概念将在后文，主要是\BookChapterRef{2}、\BookChapterRef{3}和\BookChapterRef{5}中，更深入地讨论。随后，我们概述同调镜像对称和 SYZ 背后的数学，这些内容将在\BookChapterRef{4}、\BookChapterRef{6}和\BookChapterRef{7}中展开。\BookChapterRef{8}给出 Kontsevich 最初的同调镜像猜想的精确表述，并详细处理椭圆曲线这一例子。

\section{弦理论与 \texorpdfstring{$\sigma$}{σ} 模型}\label{sec:1.1}

\SourceParagraph{ch01-1.1-001}{12}
促成这两个提案的核心物理对象，是 Polchinski 于1995年引入的 Dirichlet 膜\footnote{实际上是重新引入；有关历史，参见\BookCite{ref393,ref461}。}。这一对象在 Strominger--Yau--Zaslow 的提案中明确出现，而后来的研究也表明，它同样隐藏在 Kontsevich 的提案背后。Dirichlet 膜在物理上被定义为开弦所允许的端点，或者等价地，定义为二维共形场论中的一种边界条件。

\SourceParagraph{ch01-1.1-002}{12}
这意味着什么？虽然人们已经为共形场论和 Dirichlet 膜提出了各种有用的数学定义与解释，但就本书的目的而言，目前还没有哪一种能提供完全令人满意的起点。因此，本书的总体方法是先在直观层面解释这些物理内容，提取我们所需要的部分，然后给出可以作为更精确讨论之基础的数学定义。

\SourceParagraph{ch01-1.1-003}{12}
一般来说，弦理论描述一维弦在某个靶空间中的运动；弦在拓扑上可以是环路，即闭弦，也可以是线段，即开弦，而靶空间则是一个 Riemann 流形 $M$。为了表示一条弦随时间的运动，我们使用一个从二维 Riemann 流形 $\Sigma$ 到靶时空的映射；$\Sigma$ 称为世界面，靶时空是乘积 $M\times\mathbb{R}$，其中的 $\mathbb{R}$ 因子代表时间。

\SourceParagraph{ch01-1.1-004}{12}
要指定一个弦的量子理论，我们必须定义一个由“弦波函数”组成的 Hilbert 空间 $\mathcal{H}$，以及这个空间上的各种线性和多重线性运算。非常粗略地说，可以把 $\mathcal{H}$ 看成 $M$ 的环路空间上的泛函空间；虽然从细节上看，这个图景并不真正正确，因为波函数的支集包含不连续的环路，但它提供了一个合理的直观起点。

\SourceParagraph{ch01-1.1-005}{12}
线性运算对应于特定的世界面，或者世界面上的运算。给定一个无边界的世界面 $\Sigma$，量子理论产生一个数，即配分函数。对于 $\Sigma$ 的每一个边界，我们指定一种“边界条件”，也就是 $\mathcal{H}$ 中的一个元素，并由此得到一个数。例如，带三个边界的球面，也称“裤形曲面”，对应于 $\mathcal{H}^{\otimes3}$ 上的一个线性泛函。作用于 $\mathcal{H}$ 的其他算符，则对应于 $\Sigma$ 上度量的变化，或者其他物理可观测量。

\SourceParagraph{ch01-1.1-006}{12}
由此得到的结构，即量子场论和共形场论，与 $M$ 上的函数代数可以相互比较，并且在某种意义上是后者的推广。我们将在\BookChapterRef{3}中让读者领略这一学科的内容；不过，与几乎所有关于镜像对称的研究一样，我们的主要讨论将建立在问题的一个简化但仍十分丰富的部分之上，这就是拓扑弦理论和拓扑量子场论。

\SourceParagraph{ch01-1.1-007}{13}
在\BookChapterRef{2}中，我们将按如下思路引入拓扑弦理论。世界面与线性运算之间的对应满足“缝合关系”：一个世界面 $\Sigma$ 可以用多种方式分解成较小世界面的连通和，而对应的线性运算的各种复合必须给出相同的结果。\BookChapterRef{2}的\BookRef{fig:2.3}{图3}展示了一些简单的例子。

\SourceParagraph{ch01-1.1-008}{13}
这些缝合关系可以概括如下：

\SourceParagraph{ch01-def-1.1}{13}
\BookStatementNumber{1}
\begin{definition}\label{def:1.1}
弦理论是一个从几何范畴到线性范畴的函子。
\end{definition}

\SourceParagraph{ch01-1.1-009}{13}
我们将在\BookChapterRef{2}中讨论最简单的例子，即拓扑弦理论。在这里，所选取的几何范畴以定向的 $(d-1)$ 维流形为对象，以定向配边为态射。相应的线性范畴则可以用与之关联的有限维代数及其模来理解。

\SourceParagraph{ch01-1.1-010}{13}
同样的语言也可以用来讨论物理的量子场论，只是此时要由带度量的流形构造几何范畴。现在，Hilbert 空间 $\mathcal{H}$ 是无限维的，而态射依赖于 $\Sigma$ 上的度量。由此产生的结构只在少数情形下得到了明确描述，也就是所谓“精确可解”或“可积”的理论。由于我们需要更一般的结果，就必须讨论这些理论的物理定义。标准方法是使用对映射 $\Phi\colon\Sigma\to M\times\mathbb{R}$ 的泛函积分、相应的“微扰”图形展开，或者在某些情形下使用无限维代数的表示论。我们将在\BookChapterRef{3}中介绍这些方法。

\SourceParagraph{ch01-1.1-011}{13}
这些物理内容中的另一个重要因素是超对称。从物理上说，超对称使量子理论的性质好得多：许多原本需要通过重整化处理的麻烦发散，实际上会在费米子与玻色子的贡献之间相互抵消。超对称也是许多数学联系的核心；这一点始于 Witten 在20世纪80年代初将超对称、Morse 理论和指标理论联系起来的著名工作\BookCite{ref466,ref465,ref464,ref473}。

\SourceParagraph{ch01-1.1-012}{13}
如果假设扩展超对称，也就是说，对称代数包含若干超荷，并带有一个称为 R 对称的紧 Lie 群作用，那么理论就会受到更强的约束。这种结构是大多数与代数几何之间的联系的根源。本书主要关注的是具有“$(2,2)$”超对称的共形理论，参见\BookSectionRef{sec:3.1.4}{3.1.4}和\BookSectionRef{sec:3.3.2}{3.3.2}。在这一情形下，$M$ 必须是复 Kähler 流形。还存在其他若干情形，例如文献\BookCite{ref153}对此作了综述。

\SourceParagraph{ch01-1.1-013}{13}
将这些理论量子化时，核心的新因素是重整化群，其概要见\BookSectionRef{sec:3.2.5}{3.2.5}和\BookSectionRef{sec:3.2.6}{3.2.6}。它对靶空间 $M$ 的度量，以及在存在时对其他耦合，给出共形不变性所必需的条件。对于闭弦，在稍后将说明的意义下取最低阶，这一条件就是度量的 Ricci 平坦性；更一般地，则是超引力方程。

\SourceParagraph{ch01-1.1-014}{14}
尽管这些物理分析并不严格，但它们为大量不同的二维共形场论的存在提供了有力证据。其中一个普遍的类别，是取 $M$ 为带 Ricci 平坦度量的复 Kähler 流形。根据 Yau 对 Calabi 猜想的证明，只要 $c_1(M)=0$，这样的度量就存在。人们已经构造了大量这样的“Calabi--Yau 流形”，例如环面簇中的超曲面。

\SourceParagraph{ch01-1.1-015}{14}
由 Yau 定理还可知，Ricci 平坦度量由 $M$ 上所选取的复结构和 Kähler 类唯一确定。物理论证表明，这些 CFT 允许的形变，与复结构的无穷小变化以及复化 Kähler 类的变化一一对应。这里额外的形变对应于另一个二形式 $B$ 的形变；$B$ 必须满足调和条件，这与超引力方程一致。

\SourceParagraph{ch01-1.1-016}{14}
其他普遍的 CFT 类别包括“Landau--Ginzburg 模型”和“规范线性 $\sigma$ 模型”。它们可以看作两种数学操作的物理版本：限制到一个截面的零点集，以及对一个全纯等距变换取商。

\SourceParagraph{ch01-1.1-017}{14}
对于上述任一种模型，物理都定义了“算符谱”和“关联函数”，并给出了围绕精确可解极限作展开来计算它们的方法。对于 $\sigma$ 模型，最基本的这种极限是“大体积”极限\footnote{也称为 $\alpha'\to0$ 极限；为了读起来顺口，同时出于历史原因，还称为零斜率极限。}；在这一极限下，算符谱和关联函数化为几何不变量。一个基本例子是调和形式的代数，而以复靶空间为基础的超对称理论则可以与更精细的概念相联系，例如 Hodge 结构的变分。

\subsection{弦性修正与量子修正}\label{sec:1.1.1}

\SourceParagraph{ch01-1.1.1-001}{14}
$\sigma$ 模型的方法强调量子场论与几何之间的联系，但物理讨论中还有另一条与之相反的思路，关注的是弦理论与传统几何观念之间的差异。我们可以通过计算大体积极限的修正、使用其他更具代数色彩的共形场论方法，以及使用“半经典”论证来看到这些差异；半经典论证还将瞬子和孤子对泛函积分的额外贡献纳入考虑。许多人提出，要真正理解这些差异，最终也许需要某种新的“弦性”几何。

\SourceParagraph{ch01-1.1.1-002}{14}
先来看第一种现象，也就是大体积极限的修正。可以证明，在超对称 $\sigma$ 模型中，由重整化群分析得出的靶空间度量的共形不变性条件，实际上并不是 Ricci 平坦性，而是它的一个形变：
\begin{equation}\label{eq:1.1}
0=R_{ij}+l_s^6[R^4]_{ij}+O(l_s^8R^5).
\end{equation}
这里，$[R^4]_{ij}$ 是由 Riemann 曲率张量的四次幂构造的一个对称张量，其显式表达式见\BookCite{ref198}；$l_s$ 是一个有量纲的实形变参数，称为“弦长”。当 $l_s$ 与曲率长度相比趋近于零时，这个条件就化为 Ricci 平坦性。这些修正由量子场论的微扰理论定义，人们相信它们会延续到 $l_s^2$ 的任意阶。

\SourceParagraph{ch01-1.1.1-003}{15}
几乎所有关联函数都会得到类似的修正；可以把它们看作在大体积极限下出现的各种几何结构的形变，例如 $M$ 上调和形式的代数。然而，在如此一般的层面上，我们所知甚少；目前沿这一方向取得的几乎所有结果都来自拓扑弦理论，下面将对此讨论。

\SourceParagraph{ch01-1.1.1-004}{15}
除了弦长，弦理论中还有第二个“参数”\footnote{我们给“参数”一词加上引号，是因为可以证明，改变一个称为伸缩子的时空场，就能改变它的值。因此，从不同的 $g_s$ 值出发得到的所有理论，在物理上都是等价的。这与弦长 $l_s$ 并不是一个参数这一事实有些类似：改变 $l_s$ 的选择，总可以由整体尺度变换来补偿。}，即弦耦合，记为 $g_s$。弦耦合的定义性特点是，它控制一种展开，而其中各项来自不同亏格的世界面。例如，由涉及亏格零的球面世界面的计算得到的 Einstein 方程，可能在亏格一时获得 $g_s^2$ 阶的修正，在亏格二时获得 $g_s^4$ 阶的修正，依此类推。

\SourceParagraph{ch01-1.1.1-005}{15}
关于更高亏格的镜像对称，人们已经知道了相当多的内容；遗憾的是，本书不会涉及这个主题。也许将来《镜像对称 III》会给予它应有的篇幅。

\SourceParagraph{ch01-1.1.1-006}{15}
关于“弦性几何”的第二类信息，来自共形场论的世界面方法，或者称“非几何”方法。这些方法主要建立在 Kac--Moody 代数及其相关无限维代数的表示论之上，例如 Virasoro 代数和超 Virasoro 代数。一个著名例子是\BookSectionRef{sec:3.3.6}{3.3.6}中的“Gepner 模型”，它为某些 Calabi--Yau $\sigma$ 模型提供了独立的、原则上严格的定义。

\SourceParagraph{ch01-1.1.1-007}{15}
其中一个主题将成为我们讨论的核心，即 $N=2$ 超共形代数的理论，见\BookSectionRef{sec:3.3.3}{3.3.3}。它是镜像对称的主要物理论证的基础，见\BookSectionRef{sec:3.4.3}{3.4.3}，并将引出我们在\BookChapterRef{3}和\BookChapterRef{5}中所得的大多数具体物理结论。

\SourceParagraph{ch01-1.1.1-008}{15}
最后来看半经典方法提供的信息。这些方法纳入了具有空间延展的场构型，一般来说，它们可以分为瞬子和孤子两大类。二者都是量子场论的泛函积分定义中所用作用量泛函的非平凡临界点；这里的“非平凡”是指场构型在 $\Sigma$ 上不是常量，在 $\sigma$ 模型中也就是映射 $\Phi\colon\Sigma\to M$ 不是常值映射。这类临界点通常因拓扑上的原因而存在，但也有例外。

\SourceParagraph{ch01-1.1.1-009}{16}
瞬子是一种“集中”于底层时空 $\Sigma$ 中某一点，或者与该点关联的场构型，特别地，它也集中于时间中的某个“瞬间”。沿任意方向趋于无穷远时，它都渐近于常量。它在泛函积分的近似计算中充当“鞍点”；于是，积分被视为各个临界点贡献的总和。正如我们将在\BookChapterRef{3}中看到的那样，非平凡临界点会产生修正；与\BookEqRef{eq:1.1}{1.1}中的幂次修正不同，这种修正相对于形变参数，这里是 $l_s$，具有指数小的量级。

\SourceParagraph{ch01-1.1.1-010}{16}
在我们考虑的情形下，基本例子是取 $\Sigma\cong S^2$，并取具有非平凡 $\pi_2$ 的靶空间 $M$。这些构型称为“世界面瞬子”，会使许多关联函数得到修正。我们将在\BookSectionRef{sec:3.4}{3.4}中回顾这些修正，以及它们在镜像对称中如今已广为人知的作用。在开弦理论中，定义域 $\Sigma$ 为圆盘的映射将扮演类似的角色。

\SourceParagraph{ch01-1.1.1-011}{16}
孤子是一种与贯穿时间的线关联、但在空间上局域化的非平凡解。换言之，沿任意空间方向趋于无穷远时，场构型都趋近于常量。对于我们而言，孤子的基本例子是构成 T 对偶基础的“缠绕弦”；稍后的\BookSectionRef{sec:1.3}{1.3}会对此解释。

\SourceParagraph{ch01-1.1.1-012}{16}
当 $\Sigma$ 的维数大于二时，还可以进一步考虑在某些、但并非全部空间方向上渐近于常量的解。这些解称为“brane”，是“membrane”的简称，中文称为“膜”。如果从“时空”的角度，也就是十维的角度来考虑，那么我们即将讨论的 Dirichlet 膜就是一个例子。

\SourceParagraph{ch01-1.1.1-013}{16}
这个极简概览说明，许多物理效应都能使弦性几何与传统几何产生显著差异，但它们全都由弦长和弦耦合这两个参数控制。后续讨论中的重要参数是弦长 $l_s$：当某个几何尺度，例如曲率长度、单射半径或闭链的体积，与 $l_s$ 相比很小时，无论弦性几何究竟是什么，它都会变得重要。

\SourceParagraph{ch01-1.1.1-014}{16}
需要注意，本书有些地方，以及弦理论文献中常常采用另一种约定，以 $\alpha'=l_s^2$ 表示这个参数。

\subsection{拓扑弦理论、扭曲与镜像对称}\label{sec:1.1.2}

\SourceParagraph{ch01-1.1.2-001}{16}
要对“弦性几何”作完全一般的处理，恐怕还需要等到量子场论有了更完整、更令人满意的数学化表述。然而，这个问题中有相当重要的一部分，已经可以在现有框架中得到令人满意的理解，也就是可以放在拓扑弦理论框架内表述的那一部分。

\SourceParagraph{ch01-1.1.2-002}{17}
拓扑弦理论，或者更一般的拓扑量子场论，一个合适的基本定义是：从拓扑流形与配边所构成的范畴到线性范畴的几何函子。在这一层面上，这个主题本质上属于数学；这句话的意思是，物理方法在这里发挥不了太多作用。

\SourceParagraph{ch01-1.1.2-003}{17}
当我们能够将拓扑场论与通过泛函积分定义的物理量子场论联系起来时，物理方法就会更有价值。这种联系可以通过两种方式实现。一种是量子场论本身不依赖于 $\Sigma$ 上的度量，例如 Chern--Simons 理论\footnote{它也可能具有某种最低程度的依赖，例如 Chern--Simons 理论对标架的依赖。}。另一种是与这里有关的“上同调拓扑场论”：理论中包含一个幂零算符 $Q$，使得应力张量，也就是生成度量无穷小变化的算符，是 $Q$-恰当的。

\SourceParagraph{ch01-1.1.2-004}{17}
我们将在\BookSectionRef{sec:3.3}{3.3}中讨论 $(2,2)$ 超共形理论的这一构造，并深入介绍开弦情形所需的许多细节。这里有两种可能，即 A 扭曲与 B 扭曲，它们分别分离出物理理论中不同的、基本上相互独立的子部分。A 扭曲理论中的关联函数只依赖于复化 Kähler 模参数，见\BookSectionRef{sec:3.4.1}{3.4.1}；B 扭曲理论中的关联函数则只依赖于复结构模参数，见\BookSectionRef{sec:3.4.2}{3.4.2}。这两个理论的物理讨论十分不对称：B 模型可以完全用标准几何，也就是 Hodge 结构的变分来理解，而 A 模型中的瞬子修正，则将算符代数从经典的 de Rham 上同调环变成一个新的“量子上同调环”。

\SourceParagraph{ch01-1.1.2-005}{17}
从前面对弦性几何的讨论来看，拓扑理论之所以易于处理，是因为几乎所有幂次型的微扰修正都不再出现，只在 A 模型中留下了一系列有趣的瞬子修正。我们可以计算这些修正，例如使用泛函积分中的局部化方法，并将它们求和，得到一个显式的“弦性几何不变量”。

\SourceParagraph{ch01-1.1.2-006}{17}
正如 \BookBibText{ref242}{MS1} 详细讨论的那样，闭弦镜像对称将 Calabi--Yau 流形 $X$ 上的 A 模型，与镜像 Calabi--Yau 流形 $Y$ 上的 B 模型等同起来；通常，$Y$ 与 $X$ 之间的关系相当简单。\BookSectionRef{sec:3.4.3}{3.4.3}将概述其中对我们必不可少的部分。在很大程度上，我们可以将闭弦镜像对称的技术，如局部化、Picard--Fuchs 方程、镜像映射等，看成一个“黑箱”，在开弦讨论的特定环节调用它。

\subsection{Dirichlet 膜}\label{sec:1.1.3}

\SourceParagraph{ch01-1.1.3-001}{18}
现在可以解释为什么将 Dirichlet 膜定义为开弦所允许的端点。在开弦理论中，Hilbert 空间 $\mathcal{H}$ 大致应当是从区间到 $M$ 的映射所构成的空间上的泛函空间。当然，一个区间有两个特殊点，即起点和终点。这两个点的像各自在 $M$ 中描出一条一维轨迹，也就是“世界线”。要使开弦的定义完整，就必须规定这些端点的边界条件。

\SourceParagraph{ch01-1.1.3-002}{18}
这类边界条件最一般的定义，是用共形场论的语言表述的，它们未必具有明显的靶空间几何解释。不过，如果共形场论是以 $M$ 为靶空间的 $\sigma$ 模型，那么寻找这样的几何图景就很自然。正如\BookSectionRef{sec:3.5}{3.5}所解释的那样，这会引出如下定义。

\SourceParagraph{ch01-def-1.2}{18}
\BookStatementNumber{2}
\begin{definition}\label{def:1.2}
几何 Dirichlet 膜是一个三元组 $(L,E,\nabla_E)$：子流形 $L\subset M$ 上带有向量丛 $E$，并配备联络 $\nabla_E$。
\end{definition}

\SourceParagraph{ch01-1.1.3-003}{18}
命名时也常常加入 $L$ 的实维数：若 $p=\dim_{\mathbb{R}}L$，就称它为 Dirichlet $p$ 膜。

\SourceParagraph{ch01-1.1.3-004}{18}
从 Dirichlet 膜 $(L,E,\nabla_E)$ 伸展到 Dirichlet 膜 $(K,F,\nabla_F)$ 的一条开弦，是从区间 $I\cong[0,1]$ 到 $M$ 的映射 $X$，满足 $X(0)\in L$ 和 $X(1)\in K$。一个“开弦演化历史”是从 $\mathbb{R}$ 到开弦空间的映射；等价地，它也是从带边界的二维曲面，例如 $\Sigma\equiv I\times\mathbb{R}$，到 $M$ 的映射，使得两个边界分别嵌入 $L$ 和 $K$。

\begin{figure}[H]
\centering
\BookImage{01_Numbered_Figures/Chapter_01/Ch01_Fig01_PDFp018_Bookp008.pdf}
\caption{端点落在 D 膜上的开弦。}\label{fig:1.1}
\end{figure}

\SourceParagraph{ch01-1.1.3-005}{18}
这些开弦的量子理论，由对上述演化历史的泛函积分来定义，其权重依赖于联络 $\nabla_E$ 和 $\nabla_F$。它描述开弦态的时间演化；开弦态是 Hilbert 空间 $\mathcal{H}_{B,B'}$ 中的波函数，这个空间由两种膜的选择 $B=(L,E,\nabla_E)$ 和 $B'=(K,F,\nabla_F)$ 来标记。

\begin{figure}[H]
\centering
\BookImage{01_Numbered_Figures/Chapter_01/Ch01_Fig02_PDFp019_Bookp009.pdf}
\caption{开弦演化历史的一个例子。}\label{fig:1.2}
\end{figure}

\SourceParagraph{ch01-1.1.3-006}{19}
需要注意，不同的 Dirichlet 膜可以嵌入同一个子流形 $L$。一种表示方法，是把 Dirichlet 膜的构型指定为带重数的子流形集合；不过，也可以用\BookRef{def:1.2}{定义1.2}中对丛 $E$ 的选择来表示。因此，$N$ 个相同的膜可以通过将向量丛 $E$ 与 $\mathbb{C}^N$ 作张量积来表示，联络也由张量积得到。于是，Dirichlet 膜 $(L,E,\nabla_E)$ 的 $N$ 重副本就是三元组 $(L,E\otimes\mathbb{C}^N,\nabla_E\otimes\operatorname{id}_N)$。

\SourceParagraph{ch01-1.1.3-007}{19}
在物理中，可以给每个开弦边界标上 $\mathbb{C}^N$ 的一个基向量，来直观表示这种选择；该基向量指明了在 $N$ 个相同的膜中选取哪一个。这些标记称为“Chan--Paton 因子”，然后用来约束开弦之间的相互作用。如果把这种相互作用画成两条开弦接成一条，也就是第一条的终点接到第二条的起点，那么我们不仅要求两个端点的位置一致，也要求 Chan--Paton 因子一致。这一运算给出了“开弦代数”的直观定义。

\SourceParagraph{ch01-1.1.3-008}{19}
从数学上说，这只是说开弦代数总可以与矩阵代数作张量积，通常会得到一个非交换代数。更一般地，如果有不止一种可能的边界条件，那么与其把它看作代数，不如将它理解为由边界条件及相应开弦组成的群胚结构或范畴结构。在群胚语言中，特定的开弦是群胚的元素，复合运算仅对具有共同边界的一对开弦有定义。在范畴语言中，边界条件是对象，开弦是态射。我们将在\BookChapterRef{2}中精确表述这一思想，并在本书其余部分广泛使用它。

\SourceParagraph{ch01-1.1.3-009}{19}
为什么应当考虑非平凡的 $E$ 和 $\nabla_E$？\BookChapterRef{3}将详细说明，但说明这里可以作非平凡选择的最简单的直观论证，是援引一个一般原则：世界面作用量的任何局部形变，都应当是物理上允许的选择。开弦的端点是一个点，所以我们对点粒子的作用量能够作出的任何修改，在这里也都允许。特别地，物理中的粒子可以带有规范场的荷，例如电子对于 Maxwell 场的电荷、夸克对于色 Yang--Mills 场的荷，等等。带荷粒子的波函数便不再取复数值，而是取值于向量丛 $E$，正如上面对开弦端点的讨论一样。

\SourceParagraph{ch01-1.1.3-010}{20}
一般联络 $\nabla_E$ 的作用，是通过改变给定粒子演化历史的权重来改变泛函积分。设粒子的轨迹由映射 $\phi\colon\mathbb{R}\to M$ 定义，那么与 $M$ 上联络 $\nabla$ 关联的一个自然的轨迹泛函，就是沿该轨迹的平行移动，即从 $E|_{\phi(t_1)}$ 到 $E|_{\phi(t_2)}$ 的线性映射。物理上，此时的泛函积分被定义为对轨迹求和，并把这一平行移动包含在权重中。

\SourceParagraph{ch01-1.1.3-011}{20}
将这一图景推广到弦的最简单方法，是考虑 $l_s\to0$ 的极限。在这个极限下，只有长度趋于零、实际上等同于粒子的弦，才能满足能量有限的约束。弦的两个端点此时映到同一点，而这个点必然位于 $L\cap K$ 上。

\SourceParagraph{ch01-1.1.3-012}{20}
因此，在这一极限下，Dirichlet 膜 $(L,E,\nabla)$ 和 $(K,F,\nabla_F)$ 之间开弦的波函数，像 $L\cap K$ 上 $E^\vee\boxtimes F$ 的截面那样变换，并带有直积上的自然联络。在 $(L,E,\nabla_E)\cong(K,F,\nabla_F)$ 的特殊情形下，这就是说开弦态是 $\operatorname{End}E$ 的截面。更详细的量子化讨论将这一描述进一步细化：开弦态是分次向量丛 $\operatorname{End}E\otimes\Lambda^\bullet T^*L$ 的截面，其中次数为1的部分对应于 $\nabla_E$ 的无穷小形变。实际上，可以证明，这些开弦态正是量子场论标准意义下 $\nabla_E$ 的无穷小形变；也就是说，单条开弦是将联络 $\nabla_E$ 量子化后所得场的一个局域激发。类似地，其他开弦态是 $L$ 在 $X$ 中的法丛的截面，并以同样的方式与子流形的无穷小形变相联系。这些关系，以及它们对伸展于不同 Dirichlet 膜之间的开弦的推广，规定了如下说法的物理含义：与给定背景 $X$ 关联的那一特定集合的 Dirichlet 膜，可以从弦理论中推导出来。

\subsection{超对称、校准几何与 D 膜}\label{sec:1.1.4}

\SourceParagraph{ch01-1.1.4-001}{20}
以边界条件处理 Dirichlet 膜的物理方法，与处理“体内”量子场论的方法非常相似，下一步同样是研究重整化群。这会给出开弦所产生的场，即数据 $(M,E,\nabla)$ 的运动方程。在超引力极限下，取子流形 $M$ 在 $X$ 的度量下具有最小体积，并取联络 $\nabla$ 满足 Yang--Mills 方程，就能解这些方程。

\SourceParagraph{ch01-1.1.4-002}{21}
与 Einstein 方程一样，描述最小体积子流形的方程也是高度非线性的，其一般理论很困难。这是寻找特殊类别的解的一个动机；支持超对称的物理论证则是另一个动机。

\SourceParagraph{ch01-1.1.4-003}{21}
超对称紧化流形对应于 Ricci 平坦流形中的一个特殊类别，即允许协变常旋量的流形；类似地，Dirichlet 膜的超对称性，对应于把膜嵌入最小体积子流形中的某个特殊类别。由于物理分析建立在协变常旋量之上，这个特殊类别也应当用该旋量来定义，或者用由该旋量构造的双线性协变常形式来定义。

\SourceParagraph{ch01-1.1.4-004}{21}
导出这一类别的标准物理论证，以 Green--Schwarz 世界体作用量的 $\kappa$ 对称为基础；文献\BookCite{ref172}是一个很好的入门介绍。我们不解释这套论证，而从它的倒数第二步开始：人们发现，使度量和膜的超对称性都保持不变的那部分超对称参数 $\epsilon$，必须满足
\begin{equation}\label{eq:1.2}
\phi\equiv\operatorname{Re}\epsilon^t\Gamma\epsilon\big|_M=\operatorname{Vol}\big|_M.
\end{equation}
换言之，一个协变常形式的实部，限制到膜上时必须等于体积形式。

\SourceParagraph{ch01-1.1.4-005}{21}
由于 $\phi$ 协变常，显然有 $d\phi=0$。因此，
\[
Z(M)\equiv\int_M\phi
\]
只依赖于 $M$ 的同调类。根据不同语境，物理学家会称它为“拓扑荷”“中心荷”或“BPS 中心荷”。

\SourceParagraph{ch01-1.1.4-006}{21}
如果此外，$p$ 形式 $\phi$ 限制到任意 $p$ 维子空间 $V\subset T_xX$ 时，都不超过体积形式 $\operatorname{Vol}$，即
\begin{equation}\label{eq:1.3}
\phi|_V\leq\operatorname{Vol}|_V,
\end{equation}
那么 $\phi$ 就是 Harvey 和 Lawson 意义下的校准形式\BookCite{ref226}。这个条件可以局部检验，却蕴含一个整体结论：对于任意子流形 $M$，都有
\begin{equation}\label{eq:1.4}
\int_M\phi\leq\int_M\operatorname{Vol}.
\end{equation}
所以，中心荷的绝对值 $|Z(M)|$ 是 $\operatorname{Vol}(M)$ 的绝对下界。

\SourceParagraph{ch01-1.1.4-007}{21}
校准子流形 $M$ 就是满足\BookEqRef{eq:1.2}{1.2}的子流形，它达到上述下界，因此具有最小体积。在物理中，这类子流形通常称为“BPS 膜”；这一名称来自 Bogomol'nyi 和 Prasad--Sommerfield 对非 Abel 规范理论的磁单极子解所作的此类典型论证。

\SourceParagraph{ch01-1.1.4-008}{22}
对于 Calabi--Yau 流形 $X$，可以证明所有形式 $\omega^p/p!$ 都是校准形式；也不难证明，相应的校准子流形是复 $p$ 维的全纯子流形。此外，对任意实参数 $\theta$，$n$ 形式 $\operatorname{Re}(e^{i\theta}\Omega)$ 都是校准形式，相应的校准子流形称为特殊 Lagrange 子流形。

\SourceParagraph{ch01-1.1.4-009}{22}
上面的讨论可以推广到 $M$ 上带一般联络的情形，并由此得到 Calabi--Yau 流形 $X$ 中如下两类 BPS 膜。令 $n=\dim_{\mathbb{C}}X$，并令 $F$ 表示 $\nabla$ 的曲率二形式，其值位于 $\operatorname{End}(E)$ 中。

\SourceParagraph{ch01-1.1.4-010}{22}
第一类 BPS D 膜以校准形式 $\omega^p/p!$ 为基础，由于历史原因称为“B 型膜”。这里的 BPS 约束等价于如下三个要求：
\begin{enumerate}[label=(\arabic*)]
\item $M$ 是 $X$ 的复 $p$ 维子流形。
\item 二形式 $F$ 具有 $(1,1)$ 型，也就是说，$(E,\nabla)$ 是 $M$ 上的全纯向量丛。
\item 在超引力极限下，$F$ 满足 Hermitian Yang--Mills 方程
\[
(\omega|_M)^{p-1}\wedge F=c\cdot(\omega|_M)^p,
\]
其中 $c$ 是某个实常数。
\end{enumerate}

\SourceParagraph{ch01-1.1.4-011}{22}
计入\BookSectionRef{sec:1.1.1}{1.1.1}中的 $l_s$ 修正后，Hermitian Yang--Mills 方程形变为“MMMSL”方程\BookCite{ref347,ref330}：
\begin{enumerate}[label={$(3')$}]
\item 对某个实常数 $\phi$，$F$ 满足 $\operatorname{Im}\bigl[e^{i\phi}(\omega|_M+il_s^2F)^p\bigr]=0$。
\end{enumerate}
实际上，这一表述仍不精确；不过，进一步的修正需要更长的讨论，我们将在\BookChapterRef{3}中给出。

\SourceParagraph{ch01-1.1.4-012}{22}
第二类 BPS D 膜以校准形式 $\operatorname{Re}(e^{i\theta}\Omega)$ 为基础，称为“A 型膜”。A 膜最简单的例子是所谓特殊 Lagrange 子流形，简称 SLAG，它们满足
\begin{enumerate}[label=(\arabic*)]
\item $M$ 是 $X$ 关于 $\omega$ 的 Lagrange 子流形。
\item $F=0$，也就是说，向量丛 $E$ 是平坦的。
\item 对某个实常数 $\alpha$，有 $\operatorname{Im}(e^{i\alpha}\Omega|_M)=0$。
\end{enumerate}

\SourceParagraph{ch01-1.1.4-013}{22}
更一般地，还存在“余各向同性膜”。当 $E$ 是线丛时，这类 A 膜满足如下四个要求：
\begin{enumerate}[label=(\arabic*)]
\item $M$ 是 $X$ 关于 $\omega$ 的余各向同性子流形，即对任意 $x\in M$，$T_xM\subset T_xX$ 的辛正交补都包含于 $T_xM$。等价地，$\ker\omega_M$ 的秩等于 $M$ 在 $X$ 中的余维数；在此情形下，$\ker\omega_M$ 是 $M$ 上的可积分布。
\item 二形式 $F$ 消去 $\ker\omega_M$。
\item 令 $\mathcal{F}M$ 为向量丛 $TM/\ker\omega_M$。由前两个条件可知，$\omega_M$ 和 $F$ 下降为 $\mathcal{F}M$ 上的一对反对称形式，分别记为 $\sigma$ 和 $f$。显然，$\sigma$ 非退化。我们要求自同态 $\sigma^{-1}f\colon\mathcal{F}M\to\mathcal{F}M$ 给出 $\mathcal{F}M$ 上的复结构。
\item 令 $r$ 为 $\mathcal{F}M$ 的复维数。可以证明 $r$ 是偶数，并且 $r+n=\dim_{\mathbb{R}}M$。令 $\Omega$ 为 $K_X$ 的全纯平凡化截面。我们要求对某个实常数 $\alpha$，有 $\operatorname{Im}\bigl(e^{i\alpha}\Omega|_M\wedge F^{r/2}\bigr)=0$。
\end{enumerate}
对带有更高秩向量丛的余各向同性 A 膜，目前还没有完整的理解。

\SourceParagraph{ch01-1.1.4-014}{23}
从物理上说，还必须指定 Dirichlet 膜在时空其余 Minkowski 维度中的嵌入。最简单的可能，是取一条类时测地线，使膜在可见的四维时空中表现为粒子。只有一部分膜可以这样做，而这一部分取决于所考虑的弦理论。稍令人困惑的是，在 IIA 型理论中，B 膜是 BPS 粒子；而在 IIB 型理论中，A 膜才是 BPS 粒子。这些记号是在上述关系尚未被发现时引入的。

\subsection{弦理论与镜像对称}\label{sec:1.1.5}

\SourceParagraph{ch01-1.1.5-001}{23}
镜像对称可以用多种方式表述；对于弦理论，也许最有用的是如下表述。

\SourceParagraph{ch01-conj-1.3}{23}
\BookStatementNumber{3}
\begin{conjecture}\label{conj:1.3}
在 Calabi--Yau 三维流形 $X$ 上紧化的 IIA 型弦理论，与在镜像 Calabi--Yau 三维流形 $Y$ 上紧化的 IIB 型弦理论对偶。
\end{conjecture}

\SourceParagraph{ch01-1.1.5-002}{23}
“对偶”一词大致意味着：对于任意第一种类型的理论，都存在一个与之同构的第二种类型的理论；特别地，这两个理论的一切物理预言都相同。

\SourceParagraph{ch01-1.1.5-003}{23}
通过将各种物理可观测量相互匹配，镜像对称所有已知的数学推论都可以从这个猜想导出。我们不准备讨论它的全部影响，因为那需要深入到超出本书预期范围的物理内容。相反，我们将集中于对本书至关重要的推论，即如下猜想。

\SourceParagraph{ch01-conj-1.4}{23}
\BookStatementNumber{4}
\begin{conjecture}\label{conj:1.4}
在 $X$ 上紧化的 IIA 型理论中的 BPS D 膜集合，与在 $Y$ 上紧化的 IIB 型理论中的 BPS D 膜集合同构。
\end{conjecture}

\SourceParagraph{ch01-1.1.5-004}{23}
因为 BPS D 膜是粒子，假想的、生活在上述任一时空中的观测者都可以产生并探测它们，所以这一结论当然由主要猜想推出。提出这个猜想之后，我们可以进一步比较一对对偶理论中的 BPS Dirichlet 膜列表，以检验它。当然，这很困难。更好的方法也许是寻找某种简单构造：给定第一个列表中的一个 BPS D 膜，就能在第二个列表中产生与之对应的膜。

\SourceParagraph{ch01-1.1.5-005}{24}
另一方面，持怀疑态度的人也许会提出如下问题。我们刚才引入的所有概念，都可以用微分几何的标准语言轻易解释，而且数学界关注它们已有一段时间。如果镜像对称能够用这些概念作简单解释，那么它为什么会让数学家感到惊讶？这说明，要激发或解释镜像对称，我们需要弦理论提供某种新的东西。本书确实将用大量篇幅解释这一新输入究竟是什么；不过，实际上，有一个数学提案早于我们刚刚概述的物理，因此，先从它讲起。

\section{同调方法}\label{sec:1.2}

\SourceParagraph{ch01-1.2-001}{24}
在1994年的 ICM 报告中，Kontsevich 富有先见地提出：镜像对称可以用两个范畴之间的等价来解释，它们分别是 Calabi--Yau 流形 $X$ 上凝聚层的有界导出范畴 $D^b(X)$，以及其镜像 $Y$ 的导出 Fukaya 范畴。Fukaya 范畴的对象是 $Y$ 的 Lagrange 子流形，态射则是 Floer 上同调的元素。这两个范畴之间的导出等价，是 $X$ 与 $Y$ 的上同调同构的一个更深刻的版本；在奇数维情形下，分别对应的是 $X$ 的偶数次上同调与 $Y$ 的奇数次上同调。Kontsevich 预言，这样的等价正是镜像对称的计数预言背后的原因。

\SourceParagraph{ch01-1.2-002}{24}
这一相当抽象的提案，经过了一段时间才得到数学家和物理学家的理解与重视。它提出时，Dirichlet 膜几乎还不为人所知\footnote{当时已知的是拓扑 $\sigma$ 模型中 A 型和 B 型边界条件的定义\BookCite{ref463}，这就是 Kontsevich 的出发点。}；而被等同起来的两个范畴，也不属于数学家普遍掌握的工作知识。此外，这一提案的一大优点，同样是超越了传统微分几何和代数几何；但这也意味着，要解释其动机，就需要稍多一些弦理论知识。出于这两个原因，我们必须从更多背景材料开始讨论。

\SourceParagraph{ch01-1.2-003}{24}
因此，\BookChapterRef{3}将回顾我们所需的超共形场论内容。\BookBibText{ref242}{MS1} 已经提供了很好的入门介绍，所以我们不追求在这里给出完整论述，而是集中于以下几点。首先，CFT 与量子力学之间存在密切的类比，量子力学与数学之间的许多关系，特别是与谱几何和 Hodge 理论之间的关系，都可以简单地推广到 CFT。随后讨论 $N=2$ 超共形代数的一般理论，以及拓扑扭曲这一运算，从而与\BookChapterRef{2}中的讨论联系起来。

\SourceParagraph{ch01-1.2-004}{25}
\BookChapterRef{3}最后概述超对称边界条件。我们将从这一观点解释 A 型与 B 型 BPS 条件的起源，以及 T 对偶的物理。

\SourceParagraph{ch01-1.2-005}{25}
随后转向数学。在\BookChapterRef{4}中，我们回顾代数方面的预备知识：同调代数、凝聚层及其导出范畴，以及导出等价。

\SourceParagraph{ch01-1.2-006}{25}
箭图表示理论为\BookChapterRef{4}提供了一个理想的启发性例子。接着，我们转向同调镜像对称发挥作用的具体背景，首先回顾簇 $X$ 上凝聚层的概念。凝聚层构成一个 Abel 范畴：利用核与余核，每个态射都可以扩展为正合列。不过，层范畴在某些自然运算下，例如拉回和推出，并不具有特别好的性质。比如，拉回非局部自由层，或者对它们作张量积，都不怎么方便，而且常常造成信息丢失。簇 $X$ 的导出范畴 $D(X)$ 为此提供了解决办法。导出范畴的对象是凝聚层复形，但其态射概念比复形之间的态射更微妙，因此导出范畴乍看之下并不直观。拉回、推出、张量积和整体截面等函子，都可以扩展到导出范畴，并具有非常自然的性质。正合列的概念被正合三角形取代，为各种函子提供了长正合列的类似物。\BookChapterRef{4}将用切实易懂的方式解释这套理论，并给出充分的文献指引，供读者查阅技术细节。

\SourceParagraph{ch01-1.2-007}{25}
涉及导出范畴的一项基本运算，是 Fourier--Mukai 变换。给定簇 $X$、$Y$ 以及对象 $\mathcal{P}\in D^b(X\times Y)$，并令 $p_1$、$p_2$ 为两个投影，通过
\[
E\longmapsto p_{2*}(\mathcal{P}\otimes p_1^*E)
\]
就得到一个函子 $D^b(X)\to D^b(Y)$。这可以看作 Fourier 变换在层论中的类似物。当这种变换是范畴等价时，最为有趣，此时称为 Fourier--Mukai 函子。Mukai 最初用它证明了对偶 Abel 簇的导出范畴等价。特别地，这说明导出范畴 $D^b(X)$ 未必能确定簇 $X$。\BookChapterRef{4}将回顾 Mukai 最初的函子，以及随后对其他相对 Calabi--Yau 背景的推广，例如保典翻转和椭圆纤维化等。作为另一个说明，我们将展示如何使用 Fourier--Mukai 函子，借助简单的生成元集合 $\mathcal{O},\ldots,\mathcal{O}(-n)$，用线性代数研究射影空间 $\mathbb{P}^n$ 的导出范畴，这就是 Beilinson 的技巧。这套理论将在\BookChapterRef{4}中解释。

\SourceParagraph{ch01-1.2-008}{25}
Fourier--Mukai 变换证明有用的另一个领域，是 McKay 对应。John McKay 的一个著名观察指出，与 $\mathbb{C}^2$ 对有限子群 $G\subset\operatorname{SL}(2,\mathbb{C})$ 的商关联的 ADE 型图，只用 $G$ 的表示论就能构造出来。这在熟知的最小消解 $Y\to\mathbb{C}^2/G$ 的例外素因子与 $G$ 的非平凡不可约表示之间，建立了一一对应。如果 $G$ 是作用于 $\mathbb{C}^n$ 的 $\operatorname{SL}(n,\mathbb{C})$ 的有限子群，那么这些思想的自然推广，涉及 Nakamura 引入的 $\mathbb{C}^n$ 上 $G$ 团簇的模空间 $G\text{-}\operatorname{Hilb}(\mathbb{C}^n)$。他证明，当 $G$ 为 Abel 群时，$G\text{-}\operatorname{Hilb}(\mathbb{C}^3)$ 是 $\mathbb{C}^3/G$ 的保典消解，并猜想这一结论对所有有限子群 $G\subset\operatorname{SL}(3,\mathbb{C})$ 都成立。更一般地，只要 $Y=G\text{-}\operatorname{Hilb}(\mathbb{C}^n)\to\mathbb{C}^n/G$ 是保典消解，$Y$ 上的泛丛就确定了与 $G$ 的不可约表示 $\rho$ 一一对应的局部自由层 $\mathcal{R}_\rho$。根据 Reid 的提案，这些层应当生成 $Y$ 的 K 理论或导出范畴，就像 Beilinson 生成元生成 $\mathbb{P}^n$ 的导出范畴一样。对于有限子群 $G\subset\operatorname{SL}(3,\mathbb{C})$，Bridgeland、King 和 Reid 同时证明了 Nakamura 与 Reid 的 McKay 猜想；他们在 $Y$ 上凝聚层的有界导出范畴与 $\mathbb{C}^n$ 上 $G$ 等变凝聚层的有界导出范畴之间，建立了导出范畴等价 $\Phi\colon D(Y)\to D_G(\mathbb{C}^n)$。Craw 和 Ishii 最近的工作，对 $\mathbb{C}^3/G$ 的任意保典消解 $Y$ 建立了类似于 $\Phi$ 的导出等价，至少对于 $\operatorname{SL}(3,\mathbb{C})$ 的有限 Abel 子群是如此。我们将在\BookChapterRef{4}中解释这一系列结果背后的思想。

\SourceParagraph{ch01-1.2-009}{26}
至此，我们已经准备好深入探讨同调镜像对称的物理起源与解释，这将是\BookChapterRef{5}的内容。

\subsection{稳定性结构}\label{sec:1.2.1}

\SourceParagraph{ch01-1.2.1-001}{26}
为了说明基本要点，我们将同调镜像对称中出现的对象与 BPS 膜最初的几何描述相比较。回顾一下，A 型膜是特殊 Lagrange 流形；在同调镜像对称中，它们以 Lagrange 子流形的 Hamilton 同痕类来表示，后者是一类相关但更一般的对象。B 型膜起初看起来更相似，因为在两种描述中它们都是全纯对象，即子流形或层；不过，一旦考察联络，就会发现几何 B 型膜所带的并不只是全纯丛或层，而是更为特殊的 Hermitian Yang--Mills 联络，它满足一个额外的微分方程。

\SourceParagraph{ch01-1.2.1-002}{26}
这两类 B 型膜之间的精确关系，来自 Donaldson 定理和 Uhlenbeck--Yau 定理。它们给出了全纯丛允许不可约 Hermitian Yang--Mills 联络的充要条件：这个丛必须 $\mu$ 稳定，即在文献\BookCite{ref372}的意义下稳定。这个条件经常成立，但并不总是成立，所以几何 B 型膜组成全纯对象的一个子集。此外，在复维数至少为二时，$\mu$ 稳定性依赖于底层流形的 Kähler 类，因此几何 B 型膜的集合在 Kähler 数据的形变下并不保持不变。A 侧类似的几何考察最初由 Joyce 作出，它们表明，特殊 Lagrange 子流形组成 Lagrange 子流形的一个子集。在相应的稳定性猜想中，一个 Hamilton 同痕类是否含有特殊 Lagrange 代表，取决于一个随底层流形复结构变化的稳定性条件。因此，所预期的图景是：每个这样的类中，或者没有特殊 Lagrange 子流形，或者只有一个。

\SourceParagraph{ch01-1.2.1-003}{27}
因此，正如开弦镜像对称所要求的那样，我们得到了 A 侧与 B 侧稳定性条件之间的对应。然而，从拓扑弦理论的观点看，这颇令人惊讶，因为 A 扭曲和 B 扭曲的拓扑弦理论分别只依赖于 Kähler 模参数和复结构模参数，本来不应当对另一类模参数有任何依赖。这意味着，BPS 膜的概念不能严格地在拓扑弦理论内部定义；要作出这一定义，必须引入更多信息。

\SourceParagraph{ch01-1.2.1-004}{27}
但这个问题特别容易着手之处在于，稳定性条件只依赖于问题中几何信息的一小部分：$\mu$ 稳定性只依赖于 Kähler 类，而 Joyce 的稳定性条件只依赖于周期映射。因此，我们可以希望提出一个统一的稳定性条件，将这两种已知的几何稳定性条件都作为特例包含在内，并借此猜想：BPS 膜就是两个等价的拓扑边界条件范畴中任一范畴的稳定对象。这样一种通常称为 $\Pi$ 稳定性的条件，在 Douglas、Aspinwall 和 Bridgeland 的工作中得到了发展。

\SourceParagraph{ch01-1.2.1-005}{27}
为了发展这一图景，必须在 A 扭曲和 B 扭曲拓扑弦理论背后的共形场论中，识别出 Kontsevich 提案的组成部分。结果表明，只需要少数额外因素，它们主要来自 $N=2$ 超共形代数中 $U(1)$ 流代数部分的边界条件结构。这一部分是 Dolbeault 上同调分次背后的物理结构；计入这些物理上的选择之后，就能从原有拓扑边界条件构造分次复形，直接通向导出范畴。

\SourceParagraph{ch01-1.2.1-006}{27}
这一构造最主要的物理推论是，分次并非静态，而是动态的：改变“缺失的”模参数，例如 B 扭曲模型中的 Kähler 模参数，虽然保持了拓扑弦理论原有的全部结构，但此时新引入的分次结构却可能变化。理解这种“分次的流动”，便会相当直接地引出 $\Pi$ 稳定性的提案。

\SourceParagraph{ch01-1.2.1-007}{27}
从实际应用看，使用这一框架最有效的方式，是采用拓扑膜的 B 模型定义，并从镜像 B 模型中取得进入稳定性条件的信息，因为在这两种情况下，定义与计算都可以用标准代数几何的语言表述，不存在世界面瞬子修正。尽管一般紧 Calabi--Yau 流形上凝聚层的导出范畴还没有得到充分理解，但对于射影空间中的超曲面，可以通过限制得到许多层，从而为简单例子提供素材。此外，\BookChapterRef{4}的技术对很大一类“局部”Calabi--Yau 流形的导出范畴，给出了完整、显式的箭图描述，例如消解商奇点所得的那些流形。在闭弦一侧，周期计算由于在闭弦镜像对称中的应用，已经是一门高度发展的技术；它使我们能够展示 $\Pi$ 稳定性及其变化的许多简单例子，并用 BPS 膜的其他物理构造来检验这些例子。

\SourceParagraph{ch01-1.2.1-008}{28}
遵循 Bridgeland 的思路，进一步分析 $\Pi$ 稳定性还需要引入若干额外概念，例如广义 Harder--Narasimhan 滤过。此时，更方便的做法是把“稳定性结构”公理化，视为这些概念最一般的实现，而不再严格沿着物理讨论展开。随后可以证明，稳定性结构空间构成一个开流形，其中包含由物理例子形成的子流形。

\subsection{\texorpdfstring{$A_\infty$}{A∞} 结构的比较}\label{sec:1.2.2}

\SourceParagraph{ch01-1.2.2-001}{28}
刚才的讨论结合了 A 模型与 B 模型，在这个意义上超出了拓扑弦理论。如果询问镜像对称对拓扑弦理论有什么推论，就很自然地希望找到一个与预势类似的量；在闭弦情形下，预势编码了 Hodge 结构的变分与关联函数。

\SourceParagraph{ch01-1.2.2-002}{28}
在物理上，这个量由超势提供；可以认为，它以十分相似的方式生成开弦关联函数。用这一语言来说，镜像对称的基本计数预言，是利用适当的开弦 B 模型关联函数的级数展开，计数那些具有给定同调类、且边界位于给定特殊 Lagrange 流形上的圆盘。

\SourceParagraph{ch01-1.2.2-003}{28}
虽然这一方法已经取得成功，但它忽略了闭弦和开弦情形之间的某些关键差异，而在某些方面，这些差异比实际的计数预言更有趣。第一项差异在某种意义上很基本，却仍很重要：正如\BookSectionRef{sec:1.1.3}{1.1.3}所指出的那样，超势最好理解为非交换变量的函数。这是因为开弦关联函数对应于圆盘边界上的一组算符，而这组算符带有次序。

\SourceParagraph{ch01-1.2.2-004}{28}
更深的一项差异是：闭弦理论中出现的模空间是无阻碍的，而 Calabi--Yau 流形上向量丛的形变可能有阻碍。结果表明，这一阻碍理论恰好由超势控制：一个形变无阻碍，当且仅当超势的全部梯度都为零。因此，拓扑开弦理论的算符谱可以随着形变而变化。

\SourceParagraph{ch01-1.2.2-005}{29}
这样的结构用 $A_\infty$ 范畴描述，比用结合范畴描述更自然。它也自然地出现于圆盘边界上的关联函数构造中；事实上，人们最初就是这样发现它的，即在 Fukaya 以 Lagrange 流形与 Floer 同调为基础构造范畴时。这一结构的发展早于 Dirichlet 膜这一物理概念，并直接促成了 Kontsevich 的提案。

\SourceParagraph{ch01-1.2.2-006}{29}
回过头来看，要完全实现 Kontsevich 的提案，需要对开弦 A 模型和 B 模型作一些修改。在 A 模型一侧，Fukaya 范畴没有实现凝聚层导出范畴的三角结构；这一点可以用 Bondal 和 Kapranov 所谓的“扭曲构造”\BookCite{ref050}来补救。在 B 模型一侧，则需要稍作工作才能看到 $A_\infty$ 结构；\BookSectionRef{sec:8.2.1}{8.2.1}将遵循 Kontsevich 和 Soibelman 的工作\BookCite{ref311}来解释这一点。

\SourceParagraph{ch01-1.2.2-007}{29}
最后，在得到开弦镜像对称的精确表述后，我们将在\BookSectionRef{sec:8.4}{8.4}中详细推导 $M$ 为椭圆曲线时的基本关系，以作说明。

\section{SYZ 镜像对称与 T 对偶}\label{sec:1.3}

\SourceParagraph{ch01-1.3-001}{29}
虽然同调镜像对称的形式体系非常强大，但我们也完全可以合理地要求，对镜像对称给出其他更接近经典微分几何和代数几何的解释。这就引出了 Strominger、Yau 和 Zaslow 的提案。

\SourceParagraph{ch01-1.3-002}{29}
这一提案中的核心物理因素是 T 对偶。为了解释它，考虑以 $(M,g)$ 为靶空间的超共形 $\sigma$ 模型，把它记为
\[
\operatorname{CFT}(M,g).
\]
它可以定义为几何函子，也可以定义为一组关联函数。用物理语言来说，“对偶”是一个等价
\[
\operatorname{CFT}(M,g)\cong\operatorname{CFT}(M',g'),
\]
即使底层几何 $(M,g)$ 与 $(M',g')$ 在经典意义下并不微分同胚，这一等价仍然成立。要看到它，就必须使用\BookSectionRef{sec:1.1.1}{1.1.1}所概述的“弦性”特征。

\SourceParagraph{ch01-1.3-003}{29}
T 对偶联系两个以环面为靶空间、但具有不同度量的 CFT，其中 $M\cong M'\cong T^d$。粗略地说，这一对偶把具有长度 $L<l_s$ 的不可缩闭链的“小”靶空间，与所有这类闭链的长度均满足 $L>l_s$ 的“大”靶空间联系起来。

\SourceParagraph{ch01-1.3-004}{29}
这种关系是对偶关系的一般特征，其逻辑如下。如果所有长度尺度，包括闭链的长度、曲率长度等，都大于 $l_s$，弦理论就化为传统几何。而在传统几何中，我们知道 $(M,g)$ 与 $(M',g')$ 不同构是什么意思。对这一概念的任何修改，都必须与传统几何失效有关，而这就要求某个长度尺度“小于弦尺度”，即 $L<l_s$。

\SourceParagraph{ch01-1.3-005}{30}
为了精确表述 T 对偶，先考虑 $M=M'=S^1$。用坐标 $X\in\mathbb{R}$ 参数化它，并作等同 $X\sim X+2\pi$。考虑由 $ds^2=R^2dX^2$ 给出的 Euclid 度量 $g_R$。从它在 $\mathbb{R}^2$ 中显然的嵌入来看，实参数 $R$ 通常称为“半径”。这个流形 Ricci 平坦，所以以它为靶空间的 $\sigma$ 模型是共形场论，称为“$c=1$ 玻色子”。进一步令弦尺度 $l_s=1$。

\SourceParagraph{ch01-1.3-006}{30}
正如弦理论的基础教材\BookCite{ref395}所讨论，并且如我们将在\BookSectionRef{sec:3.2.3.6}{3.2.3.6}中证明的那样，存在完全的物理等价
\[
\operatorname{CFT}(S^1,g_R)\cong\operatorname{CFT}(S^1,g_{1/R}).
\]
因此，从弦理论的观点看，这两个靶空间无法区分。

\SourceParagraph{ch01-1.3-007}{30}
为了给出这件事的一个物理图景，姑且假定超弦理论描述我们的宇宙，因此在某种意义下，必然还存在六个额外空间维度。再假定我们有证据表明，这些额外维度在拓扑和度量上都分解为 $K_5\times S^1$。那么，询问我们宇宙中这个 $S^1$ 的半径 $R$ 是多少，就有意义。原则上，可以通过产生能量足够高的粒子，即所谓“Kaluza--Klein 模”来测量它；也可以在距离 $L\sim R$ 的尺度上测量对 Newton 引力平方反比定律的偏离。在弦理论中，T 对偶意味着 $R\geq l_s$，因为任何 $R<l_s$ 的理论都等价于另一个 $R>l_s$ 的理论。因此，两个原则上可观测的量 $R$ 与 $l_s$ 之间存在一个非平凡关系，我们可以设想通过实验检验它。

\SourceParagraph{ch01-1.3-008}{30}
回到一般讨论，考虑理论 $\operatorname{CFT}(T^d,g)$，其中 $T^d$ 是 $d$ 维环面，坐标 $X^i$ 参数化 $\mathbb{R}^d/2\pi\mathbb{Z}^d$，度量张量 $g_{ij}$ 为常量。那么存在完全的物理等价
\begin{equation}\label{eq:1.5}
\operatorname{CFT}(T^d,g)\cong\operatorname{CFT}(T^d,g^{-1}).
\end{equation}
实际上，这只是一个离散 T 对偶对称群中的一个元素；这个群由沿一维闭链的 T 对偶及大微分同胚，也就是不能与恒等映射连续相连的微分同胚，生成。完整的群同构于 $\operatorname{O}(d,d;\mathbb{Z})$\footnote{作为比较，$\operatorname{Diff}/\operatorname{Diff}_0\cong\operatorname{SL}(d,\mathbb{Z})$。}。

\SourceParagraph{ch01-1.3-009}{30}
虽然 T 对偶与传统几何十分不同，但它有一个简单的直观解释。起点是如下观察：一条弦在 $X$ 中可能的嵌入，可以按基本群 $\pi_1(X)$ 来分类。代表非平凡同伦类的弦，通常称为“缠绕态”。此外，由于弦通过在点处相互连接而发生相互作用，基于基点的环路连接所赋予 $\pi_1$ 的群结构，是有物理意义的，并且弦理论中的相互作用保持这一群结构。现在，作为 Abel 群，$\pi_1(T^d)\cong\mathbb{Z}^d$；由于显而易见的原因，它称为“缠绕数”群。

\SourceParagraph{ch01-1.3-010}{31}
当然，还可以引入另一个 $\mathbb{Z}^d$，即 $U(1)^d$ 的 Pontryagin 对偶；$T^d$ 可以看作 $U(1)^d$ 的一个仿射化。这个群的元素在物理上称为“动量”，因为它是 $T^d$ 上平移算符的本征值。相互作用同样保持这一群结构。动量和缠绕这两种群结构，可以概括为：完整闭弦代数包含群代数 $\mathbb{C}[\mathbb{Z}^d]\oplus\mathbb{C}[\mathbb{Z}^d]$。

\SourceParagraph{ch01-1.3-011}{31}
本质上，T 对偶的要点是：如果在足够小的靶空间上将弦量子化，动量与缠绕的角色就会互换。\BookSectionRef{sec:3.2.3.6}{3.2.3.6}将用一个简短的泛函积分论证说明这一点。不过，引入一些初等谱几何，就能看出主要思想。除了刚才讨论的代数结构，共形场论的另一个不变量是其 Hamilton 算符 $H$ 的谱；选定时间单位，并除去与态无关的真空能量平移后，这个算符由 Virasoro 算符 $L_0+\overline{L}_0$ 给出。这个 Hamilton 算符可以看作 $X$ 上函数的标准 Laplace 算符 $\Delta_g$ 的类似物。不难看出，对于带有上述度量 $g$ 的 $T^d$，其谱为
\begin{equation}\label{eq:1.6}
\operatorname{Spec}\Delta_g=\left\{\sum_{i,j=1}^d g^{ij}p_i p_j\;;\;p\in\mathbb{Z}^d\right\}.
\end{equation}
另一方面，正如直观上可以预料的那样，缠绕弦的能量是其长度的函数。在这个环面上，缠绕数为 $w\in\mathbb{Z}^d$ 的测地线，其长度平方为
\begin{equation}\label{eq:1.7}
L^2=(2\pi)^2\sum_{i,j=1}^d g_{ij}w^iw^j.
\end{equation}
现在，只需从弦理论引入一项事实：总 Hamilton 算符同时包含这两项，
\[
H=\Delta_g+\frac{L^2}{(2\pi)^2}+\cdots,
\]
其中额外的 $\cdots$ 项表示弦的激发模，也就是“振子”模的能量。于是，取逆 $g\to g^{-1}$，同时交换 $p\leftrightarrow w$，就使 $H$ 的谱保持不变。这就是 T 对偶。

\SourceParagraph{ch01-1.3-012}{31}
上述讨论可以简单地推广到环面上存在非零 B 场且满足 $dB=0$ 的情形。此时，$B$ 是常量反对称张量，可以用矩阵 $g+B$ 标记 CFT。基本的 T 对偶关系变为
\[
\operatorname{CFT}(T^d,g+B)\cong\operatorname{CFT}\bigl(T^d,(g+B)^{-1}\bigr).
\]

\SourceParagraph{ch01-1.3-013}{32}
另一种推广要微妙得多，是对族作 T 对偶，或者说逐纤维作 T 对偶。同样的论证仍然适用；如果纤维度量的变化长度尺度远大于 $l_s$，且曲率长度也远大于 $l_s$，那么在这一极限下，论证就会变得精确。物理中有时称它为“绝热极限”。

\SourceParagraph{ch01-1.3-014}{32}
虽然这个假设非常苛刻，但还有更具启发性的物理论证，表明 T 对偶应当在更一般的情形下成立；我们所讨论的关系会得到与纤维度量的曲率 $l_s^2R$ 和导数 $l_s\partial$ 成比例的修正，这些修正既来自微扰理论，也来自世界面瞬子。这些修正尚未得到充分研究，令人遗憾，因为正如下文将讨论的那样，它们很可能有助于阐明精确表述 SYZ 猜想时遇到的微妙问题。

\subsection{T 对偶与 Dirichlet 膜}\label{sec:1.3.1}

\SourceParagraph{ch01-1.3.1-001}{32}
刚才的讨论针对闭弦。显然，从区间到流形的映射并不由 $\pi_1$ 分类，事实上也没有类似的选择。那么，T 对偶关系如何能对开弦成立呢？

\SourceParagraph{ch01-1.3.1-002}{32}
正是这个问题导致了 Dirichlet 膜最初的发现\BookCite{ref106}。假设从能够在 $T^d$ 中任意自由传播的开弦出发；用现代语言说，就是有一个包裹 $T^d$ 的 Dirichlet $d$ 膜。虽然此时没有缠绕数，但这种弦态是 $T^d$ 上的函数，因此仍可应用 Pontryagin 对偶，得到开弦代数包含 $\mathbb{C}[\mathbb{Z}^d]$ 的结论。用物理语言说，存在一个 $d$ 维的守恒动量。此外，开弦的 Hamilton 算符仍包含一个类似 $T^d$ 上 Laplace 算符的部分，它的谱仍然由\BookEqRef{eq:1.6}{1.6}给出。

\SourceParagraph{ch01-1.3.1-003}{32}
如果作取逆变换 $g\to g^{-1}$，显然，恢复原谱的最简单办法，是找到一个新的开弦部分，使其谱再次由\BookEqRef{eq:1.7}{1.7}中的 $L^2/(2\pi)^2$ 的可能取值组成。对于包裹 $T^d$ 的 $d$ 膜上的开弦，这并不成立；不过，如果强制开弦的两个端点重合，就可能成立，因为满足这一条件的测地线的最小长度平方，同样由\BookEqRef{eq:1.7}{1.7}给出。此外，$\pi_1$ 通常用带基点的环路定义，当然在这里有无基点都不影响结果，所以，强制在特定点 $p$ 开始并结束的开弦所组成的部分，仍将包含代数 $\mathbb{C}[\mathbb{Z}^d]$。

\SourceParagraph{ch01-1.3.1-004}{32}
因此，一个简单的提案是：包含 Dirichlet $d$ 膜的理论 $\operatorname{CFT}(T^d,g)$，其 T 对偶是包含 Dirichlet $0$ 膜的理论 $\operatorname{CFT}(T^d,g^{-1})$。这一提案将在\BookSectionRef{sec:3.5.4}{3.5.4}中用泛函积分论证来证明。这样的 $0$ 膜由一个零维子流形的选择来定义，也就是一个点 $p\in X$，其开弦必须在 $p$ 处开始或结束。

\SourceParagraph{ch01-1.3.1-005}{33}
这恢复了原谱，但现在必须追问：点 $p$ 的选择，在原来的 $d$ 膜理论中对应于什么？这个问题有一个优美的答案：为了完整地指定 $d$ 膜，还必须指定一个丛 $E$ 和一个联络 $\nabla$。$0$ 膜的 T 对偶，是带平凡丛 $E$ 和平坦联络的 $d$ 膜，这既合乎情理，也确实正确。不过，$T^d$ 上平坦联络的模空间本身就是一个环面，记为 $\widetilde{T}^d$。由此可以检验一个稍不平凡的结论：将带度量 $g$ 的 $T^d$ 上联络空间的自然度量限制到平坦联络，所得到的 $\widetilde{T}^d$ 上的自然度量，正是平坦度量 $g^{-1}$。虽然联络空间度量的整体尺度尚未确定，但弦理论确定这一关系恰好是\BookEqRef{eq:1.5}{1.5}。因此，所提出的这一对 T 对偶 Dirichlet 膜理论的模空间，作为度量空间是相同的。

\subsection{镜像对称与特殊 Lagrange 纤维化}\label{sec:1.3.2}

\SourceParagraph{ch01-1.3.2-001}{33}
现在可以解释 Strominger--Yau--Zaslow 提案\BookCite{ref433}。考虑一对由镜像对称联系起来的紧 Calabi--Yau 三维流形 $X$ 和 $Y$。根据前面的讨论，$X$ 上的 BPS A 膜集合同构于 $Y$ 上的 BPS B 膜集合，而 $X$ 上的 BPS B 膜集合同构于 $Y$ 上的 BPS A 膜集合。

\SourceParagraph{ch01-1.3.2-002}{33}
$X$ 上最简单的 BPS B 膜是点。对于 $X$ 的所有复结构，包括非代数的复结构，这些膜都存在，其模空间就是 $X$ 本身。试着确定它们对应于 $Y$ 上哪些 BPS A 膜。前面描述的 BPS A 膜条件意味着，A 膜可以具有实维数5或3。但是，五维 A 膜的存在条件对辛结构十分敏感；例如，将辛形式乘以一个常数因子来重标度，一般会消除这类 A 膜，因为线丛的曲率二形式具有量子化的周期。相比之下，将辛形式乘以任意常数因子后，特殊 Lagrange 子流形仍然是特殊 Lagrange 子流形。这与 $X$ 上点的性质吻合。因此，$X$ 上一个点的镜像必须是一个三维 A 膜 $(N,E,\nabla)$。在这一情形下，BPS 条件意味着 $F=0$，这里为简单起见假设 $B=0$，并且 $N$ 是 $Y$ 的特殊 Lagrange 子流形。因此可以得出结论：$Y$ 上存在一个以 $X$ 的点为参数的特殊 Lagrange 子流形族。而且，这个族就是将单个特殊 Lagrange 子流形 $N$ 视为 $Y$ 上 BPS A 膜时所得的模空间。

\SourceParagraph{ch01-1.3.2-003}{33}
对于 $Y$ 上这个特殊的特殊 Lagrange 子流形族，还能说些什么？根据 McLean 定理，见\BookSectionRef{sec:6.1.1}{6.1.1}，特殊 Lagrange 子流形 $N$ 的模空间局部光滑，维数为 $b_1(N)$。一个 BPS A 膜，是带 Hermitian 线丛 $E$ 及其平坦联络 $\nabla$ 的特殊 Lagrange 子流形。在固定的 $N$ 上，平坦联络的模空间是实维数为 $b_1(N)$ 的环面。因此，$(N,E,\nabla)$ 的模空间总维数为 $2b_1(N)$。这个模空间是 $X$，所以必有 $b_1(N)=3$。于是 $X$ 由三维环面纤维化。交换 $Y$ 与 $X$ 的角色，也得到 $Y$ 由三维环面纤维化的结论。

\SourceParagraph{ch01-1.3.2-004}{34}
由此得到的 $X$ 和 $Y$ 的 $T^3$ 纤维化，不可能处处光滑，否则 $X$ 和 $Y$ 的 Euler 示性数都会为零。奇点出现在形变后的特殊 Lagrange 子流形不再光滑的时候。

\SourceParagraph{ch01-1.3.2-005}{34}
接下来希望说明，这两个 $T^3$ 纤维化的光滑纤维本身也是特殊 Lagrange 子流形。设想 $X$ 的 $T^3$ 纤维化的纤维上的诱导度量是平坦的。这个假设并不现实，但可以希望，在大体积极限下，远离奇异纤维时它是一个好的近似。那么，就可以在非奇异纤维上作 T 对偶，得到一个非紧 Calabi--Yau 流形 $X'$，它与 $X$ 的一个开部分互为镜像。T 对偶把位于纤维 $M$ 中的点 $p\in X$ 映为形如 $(M',E,\nabla)$ 的 D 膜，其中三维环面 $M'$ 与 $M$ 对偶，而平坦联络 $\nabla$ 由 $p$ 在 $M$ 中的位置确定。让 $p$ 在 $X$ 上移动，既会使 $M'$ 形变，也会使平坦联络 $\nabla$ 形变。这强烈暗示，$X'$ 是 $Y$ 的一个开部分，而 $(M',E,\nabla)$ 就是与 $p$ 互为镜像的 A 膜 $N$。于是，$X$ 和 $Y$ 的 $T^3$ 纤维化近似互为 T 对偶。此外，如果考虑 $N=M'$ 上的一点，那么它的 T 对偶就是 $M$。但 $Y$ 上任一点都是 BPS B 膜，所以它的 T 对偶 $M$ 必须是 $X$ 上的 BPS A 膜。因此，$X$ 的 $T^3$ 纤维化的每个非奇异纤维都是特殊 Lagrange 子流形。交换 $X$ 和 $Y$ 的角色，也可得出 $Y$ 的 $T^3$ 纤维化的非奇异纤维同样是特殊 Lagrange 子流形。

\SourceParagraph{ch01-1.3.2-006}{34}
我们将这些讨论概括为如下猜想。

\SourceParagraph{ch01-conj-1.5}{34}
\BookStatementNumber{5}
\begin{conjecture}[Strominger--Yau--Zaslow]\label{conj:1.5}
对于任意一对互为镜像的紧、单连通 Calabi--Yau 三维流形 $X$ 与 $Y$，应当存在 $X$ 和 $Y$ 的 $T^3$ 纤维化，并具有如下两条性质：
\begin{itemize}
\item 非奇异纤维是特殊 Lagrange 子流形。
\item 对 $X$ 取大体积极限，并对 $Y$ 取相应的“大复结构”极限后，这两个纤维化互为 T 对偶。
\end{itemize}
\end{conjecture}

\SourceParagraph{ch01-1.3.2-007}{34}
由于物理论证建立在亏格零的 CFT 之上，所以没有明显理由将其限制在六个实维度和 $\hat{c}=3$ 的情形。因此，还可以进一步猜想，SYZ 提案不仅适用于三维流形，也适用于任意维数 $d$ 的 Calabi--Yau 镜像对；此时，它预言的是 $T^d$ 纤维化结构。

\subsection{SYZ 猜想的数学}\label{sec:1.3.3}

\SourceParagraph{ch01-1.3.3-001}{34}
正如已经看到的那样，物理给出了相当简单的镜像对称图景：对于非奇异纤维，只需对环面作 T 对偶，就能得到镜像。另一方面，SYZ 的工作没有解释如何处理奇异纤维，后续物理研究也没有真正填补这一空缺。

\SourceParagraph{ch01-1.3.3-002}{35}
不过，自1996年最初的论文发表以来，数学界在理解这一猜想方面取得了许多进展。Hitchin 的重要工作描述了特殊 Lagrange 纤维化的底空间上出现的自然结构。设 $f\colon X\to B$ 是\BookRef{conj:1.5}{猜想1.5}所预期的特殊 Lagrange 纤维化。如果
\[
B_0=\{b\in B\mid f^{-1}(b)\text{是非奇异环面}\},
\]
那么 $B_0$ 带有一个所谓 Hessian 流形结构，或者称仿射 Kähler 流形结构。对于对偶纤维化，这些结构应当由一个自然的对偶过程联系起来，而这个过程本质上只是 Legendre 变换。

\SourceParagraph{ch01-1.3.3-003}{35}
我们将在\BookSectionRef{sec:6.2}{6.2}中详细描述这些结构，并据此完整地呈现半平坦镜像对称中出现的结构。在这一情形下，$B=B_0$，且 $X$ 上的 Ricci 平坦度量限制到 $f$ 的每个纤维上都是平坦度量。此时，$X$ 的全部信息都可以从 $B$ 上的结构恢复出来，人们也由此开始意识到，镜像对称中的关键对象并不是 Calabi--Yau 流形本身，而是仿射 Kähler 底空间。我们将对此作深入探讨。我们能够看到镜像对称所许诺的复结构模参数与 Kähler 模参数之间的镜像同构，并且还将从更深的层面描述镜像对称：它把 $X$ 上某些 Lagrange 子流形与 $Y$ 上的向量丛或凝聚层相互交换。这说明了 SYZ 猜想与同调镜像对称猜想之间的联系。这部分内容将在\BookSectionRef{sec:6.3}{6.3}中介绍。

\SourceParagraph{ch01-1.3.3-004}{35}
为了更充分地探讨 SYZ 猜想，我们将采用几种方法。首先，作为\BookChapterRef{6}中对 SYZ 猜想初步讨论的结尾，我们将说明如何构造能够紧化为非平凡 Calabi--Yau 流形的环面纤维化，例如 $\mathbb{P}^4$ 中的五次超曲面。这需要理解奇异纤维，\BookChapterRef{6}将介绍这些内容。在这一情形下，我们能够构造拓扑环面纤维化，其中环面在余维二的集合 $\Delta\subset B$ 上退化，这个集合称为判别集。对于五次超曲面，这个判别集是一张三价图。不过，这是纯拓扑的方法，并未处理猜想的度量方面。

\SourceParagraph{ch01-1.3.3-005}{35}
度量问题将在\BookChapterRef{7}中处理，届时我们将离开易于处理的半平坦情形。我们首先从若干观点出发，描述并构造 Ricci 平坦度量及其他特殊度量的例子，其中会用到 Hitchin 最近的工作。

\SourceParagraph{ch01-1.3.3-006}{35}
对于紧 Calabi--Yau 流形，Yau 的一个深刻定理保证，对流形上的每个 Kähler 类，都存在一个 Kähler Ricci 平坦度量。然而，人们还不知道任何一个显式、非平凡的例子。另一方面，在非紧情形下，如果存在对称性，往往可以将支配 Ricci 平坦性的复杂偏微分方程，化为常微分方程。在这种情形下，一些一般的拟设可供使用，例如 Gibbons--Hawking 拟设。\BookChapterRef{7}将讨论若干选取的例子。

\SourceParagraph{ch01-1.3.3-007}{36}
\BookChapterRef{7}也将介绍 Joyce 关于特殊 Lagrange 纤维化的一些重要的新发现。\BookChapterRef{6}中讨论的半平坦和拓扑 SYZ 图景，与最初的 SYZ 图景很吻合；但 Joyce 发展的一般特殊 Lagrange 纤维化图景，却与之不符。特别地，他打破了特殊 Lagrange 纤维化必然是 $C^\infty$ 的希望；因此，我们不再预期特殊 Lagrange 纤维化的判别集必然是一个性质良好的余维二对象。于是，\BookChapterRef{6}中为五次超曲面建立的、以图为判别集的拓扑图景，似乎只是对假想的特殊 Lagrange 纤维化的一个近似；如今，我们预期其判别集会是这张图的某种增厚。此外，Joyce 提供的证据暗示，对偶化过程实际上应当改变这个增厚的判别集。因此，最初的 SYZ 强形式，也就是纤维之间存在精确对偶的形式，不可能成立，因为对偶纤维化中可能出现 $f^{-1}(b)$ 非奇异而 $\check{f}^{-1}(b)$ 奇异的情形。于是，我们不得不修改 SYZ 猜想的度量版本。

\SourceParagraph{ch01-1.3.3-008}{36}
Kontsevich 和 Soibelman 以及 Gross 和 Wilson 的工作发展了这一修订。基本思想是，我们只预期 SYZ 在所谓大复结构极限点附近成立。实质上，这意味着给定一个 Calabi--Yau 流形的退化族 $\pi\colon\mathcal{X}\to S$，并有一点 $0\in S$，使 $\pi^{-1}(0)$ 是一个极度奇异的 Calabi--Yau 簇；“极度奇异”在这里可以精确定义。对于 $S$ 中接近0的点 $t$，我们预期纤维 $X_t$ 带有特殊 Lagrange 纤维化，虽然它可能具有相当差的性质，而且可能只在 $X_t$ 的某个开子集上存在。不过，当 $t\to0$ 时，我们希望这个特殊 Lagrange 纤维化的性质会改善，并“趋向”某个性质良好的极限。\BookChapterRef{7}将精确表述这一点。现有证据暗示，这与 $t\to0$ 时 $X_t$ 上 Ricci 平坦度量的整体行为密切相关。特别地，我们可以精确定义度量空间极限的概念，并猜想，当 $t\to0$ 时，$X_t$ 实际上发生塌缩：特殊 Lagrange 环面缩为点，所得极限是维数为 $X_t$ 一半的流形。这个极限流形应当恰好与特殊 Lagrange 纤维化的底空间一致。这使我们能够提出 SYZ 猜想的极限形式；目前看来，它是这一猜想最可行的微分几何形式。
